\subsection{Hipertor de dos complejos de módulos}
\label{subsection:III.6.3}

\begin{env}[6.3.1]
\label{III.6.3.1}
Sean $A$ un anillo conmutativo y $P_\bullet$, $Q_\bullet$ dos complejos de
$A$-módulos cuyos diferenciales son de grado $-1$; sea
$L_{\bullet\bullet}$ (resp. $M_{\bullet\bullet}$) una resolución proyectiva
de Cartan-Eilenberg de $P_\bullet$ (resp. $Q_\bullet$)
\sref[0]{0.11.6.1};
$L_{\bullet\bullet}\otimes_A M_{\bullet\bullet}$ es entonces (para la suma
de los primeros grados y la suma de los segundos grados) un bicomplejo (con
diferenciales de grado $-1$), cuya homología
$H_\bullet(L_{\bullet\bullet}\otimes_A M_{\bullet\bullet})$ no depende de las
resoluciones de Cartan-Eilenberg $L_{\bullet\bullet}$,
$M_{\bullet\bullet}$ elegidas y es, por definición, la hiperhomología del
bifuntor $P_\bullet\otimes_A Q_\bullet$ en $P_\bullet$ y $Q_\bullet$
\sref[0]{0.11.6.5}. Pondremos, por definición,
\[
  \mathbf{Tor}_n^A(P_\bullet,Q_\bullet)
  =H_n(L_{\bullet\bullet}\otimes_A M_{\bullet\bullet})
  \tag{6.3.1.1}\label{III.6.3.1.1}
\]

\oldpage[III]{142}

y diremos que este $A$-módulo es el hipertor de índice $n$ de los dos
complejos $P_\bullet$, $Q_\bullet$. Se sabe que, en la categoría de los
complejos de $A$-módulos acotados inferiormente, los
$\mathbf{Tor}_n^A(P_\bullet,Q_\bullet)$ forman un bifuntor homológico en
$P_\bullet$, $Q_\bullet$ \sref[0]{0.11.6.5}. Además:
\end{env}

\begin{proposition}[6.3.2]
\label{III.6.3.2}
El bifuntor $\mathbf{Tor}_\bullet^A(P_\bullet,Q_\bullet)$ es el término
límite común de dos bifuntores espectrales
$'E(P_\bullet,Q_\bullet)$, $''E(P_\bullet,Q_\bullet)$, cuyos términos $E_2$
son
\begin{align}
  'E^2_{pq}
    &=H_p\bigl(\operatorname{Tor}_q^A(P_\bullet,Q_\bullet)\bigr),
    \tag{6.3.2.1}\label{III.6.3.2.1}\\
  ''E^2_{pq}
    &=\bigoplus_{q'+q''=q}
      \operatorname{Tor}_p^A\bigl(H_{q'}(P_\bullet),H_{q''}(Q_\bullet)\bigr).
    \tag{6.3.2.2}\label{III.6.3.2.2}
\end{align}
donde, en \sref{III.6.3.2.1},
$\operatorname{Tor}_q^A(P_\bullet,Q_\bullet)$ designa el
bicomplejo formado por los $A$-módulos $\operatorname{Tor}_q^A(P_i,Q_j)$. La
sucesión espectral \sref{III.6.3.2.2} es siempre regular; si $P_\bullet$ y
$Q_\bullet$
están acotados inferiormente, o si $A$ es de dimensión cohomológica finita,
las dos sucesiones espectrales \sref{III.6.3.2.1} y
\sref{III.6.3.2.2} son birregulares.
\end{proposition}

\begin{proof}
Esto resulta de \sref[0]{0.11.6.5}, pues cuando $A$ es de dimensión
cohomológica
finita $n$, todo $A$-módulo admite una resolución proyectiva de longitud $n$
\emph{(M, VI, 2.1)}.
\end{proof}

\begin{corollary}[6.3.3]
\label{III.6.3.3}
Sean $P'_\bullet$, $Q'_\bullet$ dos complejos de $A$-módulos, y
$u:P_\bullet\to P'_\bullet$, $v:Q_\bullet\to Q'_\bullet$ dos homomorfismos
de complejos. Si los homomorfismos
$H_\bullet(u):H_\bullet(P_\bullet)\to H_\bullet(P'_\bullet)$,
$H_\bullet(v):H_\bullet(Q_\bullet)\to H_\bullet(Q'_\bullet)$ deducidos
respectivamente de $u$ y $v$ son biyectivos, entonces el homomorfismo
$\mathbf{Tor}_\bullet^A(P_\bullet,Q_\bullet)\to
\mathbf{Tor}_\bullet^A(P'_\bullet,Q'_\bullet)$ deducido de $u$ y $v$ es
biyectivo.
\end{corollary}

\begin{proof}
En efecto, el homomorfismo de sucesiones espectrales
$''E(P_\bullet,Q_\bullet)\to''E(P'_\bullet,Q'_\bullet)$ deducido de $u$ y $v$
es entonces un isomorfismo para los términos $E_2$, y la conclusión resulta
de que estas sucesiones son regulares en virtud de \sref{III.6.3.2} y
\sref[0]{0.11.1.5}.
\end{proof}

\begin{proposition}[6.3.4]
\label{III.6.3.4}
Sean $P_\bullet$, $Q_\bullet$ dos complejos de $A$-módulos acotados
inferiormente. Sea $L_{\bullet\bullet}$ (resp. $M_{\bullet\bullet}$) un
bicomplejo formado por $A$-módulos planos, tal que, para todo $i$,
$L_{i,\bullet}$ (resp. $M_{i,\bullet}$) sea una resolución de $P_i$
(resp. $Q_i$). Se tienen entonces isomorfismos canónicos
\[
  \mathbf{Tor}_\bullet^A(P_\bullet,Q_\bullet)
  \xrightarrow{\sim}H_\bullet(L_{\bullet\bullet}\otimes_AQ_\bullet)
  \xrightarrow{\sim}H_\bullet(P_\bullet\otimes_AM_{\bullet\bullet})
  \xrightarrow{\sim}H_\bullet(L_{\bullet\bullet}\otimes_AM_{\bullet\bullet})
  \tag{6.3.4.1}\label{III.6.3.4.1}
\]
\end{proposition}

\begin{proof}
Esto resulta de \sref[0]{0.11.6.5}, partes (ii) y (iii), y de la definición
de los
$A$-módulos planos.
\end{proof}

\begin{namedtheorem}[6.3.5]{Observaciones}
\label{III.6.3.5}
(i) Con las notaciones de \sref{III.6.3.1}, los bicomplejos
$L_{\bullet\bullet}\otimes_A M_{\bullet\bullet}$ y
$M_{\bullet\bullet}\otimes_A L_{\bullet\bullet}$ son canónicamente
isomorfos, de donde se obtiene un isomorfismo canónico
$\mathbf{Tor}_\bullet^A(P_\bullet,Q_\bullet)\xrightarrow{\sim}
\mathbf{Tor}_\bullet^A(Q_\bullet,P_\bullet)$.

(ii) Si $F$ y $G$ son dos $A$-módulos, y $P_\bullet$ y $Q_\bullet$ son los
complejos de $A$-módulos reducidos a $F$ y $G$, respectivamente, en grado 0 y
nulos en los demás grados, entonces dos resoluciones proyectivas $L_\bullet$,
$M_\bullet$ de $F$ y $G$, respectivamente, pueden considerarse como
resoluciones de Cartan-Eilenberg de $P_\bullet$ y $Q_\bullet$ completándolas
con ceros. Por consiguiente, en este caso se tiene
$\mathbf{Tor}_\bullet^A(P_\bullet,Q_\bullet)
=\operatorname{Tor}_\bullet^A(F,G)$.
\end{namedtheorem}

\begin{proposition}[6.3.6]
\label{III.6.3.6}
Sean $(P_\bullet^\lambda)$, $(Q_\bullet^\mu)$ dos sistemas inductivos
filtrantes de complejos de $A$-módulos; se tiene un isomorfismo canónico
\[
  \varinjlim_{\lambda,\mu}
  \mathbf{Tor}_\bullet^A(P_\bullet^\lambda,Q_\bullet^\mu)
  \xrightarrow{\sim}
  \mathbf{Tor}_\bullet^A
  \left(\varinjlim_\lambda P_\bullet^\lambda,
        \varinjlim_\mu Q_\bullet^\mu\right).
  \tag{6.3.6.1}\label{III.6.3.6.1}
\]
\end{proposition}

\begin{proof}
Pongamos $P_\bullet=\varinjlim_\lambda P_\bullet^\lambda$,
$Q_\bullet=\varinjlim_\mu Q_\bullet^\mu$; por funtorialidad, está claro que
los $\mathbf{Tor}_\bullet^A(P_\bullet^\lambda,Q_\bullet^\mu)$ forman un
sistema inductivo y que las aplicaciones
$\mathbf{Tor}_\bullet^A(P_\bullet^\lambda,Q_\bullet^\mu)\to
\mathbf{Tor}_\bullet^A(P_\bullet,Q_\bullet)$ deducidas

\oldpage[III]{143}

de las aplicaciones canónicas $P_\bullet^\lambda\to P_\bullet$,
$Q_\bullet^\mu\to Q_\bullet$, forman un sistema inductivo de homomorfismos,
de donde se obtiene un homomorfismo canónico \sref{III.6.3.6.1} y, más
generalmente,
un homomorfismo canónico
\[
  \varinjlim_{\lambda,\mu}''E(P_\bullet^\lambda,Q_\bullet^\mu)
  \longrightarrow ''E(P_\bullet,Q_\bullet)
\]
del cual \sref{III.6.3.6.1} es el homomorfismo de los términos límite.
Además, la sucesión espectral $''E(P_\bullet,Q_\bullet)$ es regular
\sref{III.6.3.2}, y lo mismo
ocurre con la sucesión espectral
$\varinjlim_{\lambda,\mu}''E(P_\bullet^\lambda,Q_\bullet^\mu)$, como resulta
de las definiciones \sref[0]{0.11.1.7} y de la demostración de
\sref[0]{0.11.3.3}; para demostrar que \sref{III.6.3.6.1} es biyectivo basta,
pues, por \sref[0]{0.11.1.5}, probar que el
homomorfismo
\[
  \varinjlim_{\lambda,\mu}''E(P_\bullet^\lambda,Q_\bullet^\mu)
  \longrightarrow ''E(P_\bullet,Q_\bullet)
  \tag{6.3.6.2}\label{III.6.3.6.2}
\]
es biyectivo para los términos $E^2$. Como el funtor $H_\bullet$ conmuta con
el límite inductivo de complejos de módulos, finalmente se reduce todo a
probar que, para dos sistemas inductivos filtrantes $(F^\lambda)$,
$(G^\mu)$ de $A$-módulos, el homomorfismo canónico
\[
  \varinjlim_{\lambda,\mu}\operatorname{Tor}_\bullet^A(F^\lambda,G^\mu)
  \longrightarrow
  \operatorname{Tor}_\bullet^A
  \left(\varinjlim_\lambda F^\lambda,
        \varinjlim_\mu G^\mu\right)
\]
es biyectivo. Para ello, consideremos para cada $F^\lambda$ la resolución
libre canónica
\[
  L_\bullet^\lambda:\quad
  \cdots\longrightarrow L_{i+1}^\lambda
  \longrightarrow L_i^\lambda\longrightarrow\cdots
  \longrightarrow L_1^\lambda\longrightarrow L_0^\lambda\longrightarrow0
\]
donde $L_0^\lambda$ es el $A$-módulo de las combinaciones lineales formales
de elementos de $F^\lambda$, y $L_{i+1}^\lambda$ es el $A$-módulo de las
combinaciones lineales formales de elementos de
$\operatorname{Ker}(L_i^\lambda\to L_{i-1}^\lambda)$; se comprueba
inmediatamente que los $L_\bullet^\lambda$ forman un sistema inductivo de
complejos, y si se pone
$F=\varinjlim_\lambda F^\lambda$, $L_i=\varinjlim_\lambda L_i^\lambda$, los
$L_i$ forman una resolución $L_\bullet$ de $F$, pues el funtor
$\varinjlim$ es exacto; además, los $L_i$, límites inductivos de
$A$-módulos libres, son planos \sref[0]{0.6.1.2}. Se considera de
igual modo, para cada $\mu$, la resolución libre canónica $M_\bullet^\mu$ de
$G^\mu$, y $M_\bullet=\varinjlim_\mu M_\bullet^\mu$ es una resolución plana
de $G=\varinjlim_\mu G^\mu$. Se tiene entonces
$\operatorname{Tor}_\bullet^A(\varinjlim_\lambda F^\lambda,
\varinjlim_\mu G^\mu)=H_\bullet(L_\bullet\otimes_A M_\bullet)$ en virtud de
\sref{III.6.3.5} y \sref{III.6.3.4}; pero
\[
  H_\bullet(L_\bullet\otimes_A M_\bullet)
  =\varinjlim_{\lambda,\mu}
   H_\bullet(L_\bullet^\lambda\otimes_A M_\bullet^\mu)
\]
puesto que $H_\bullet$ conmuta con los límites inductivos de complejos de
módulos; como $H_\bullet(L_\bullet^\lambda\otimes_A M_\bullet^\mu)
=\operatorname{Tor}_\bullet^A(F^\lambda,G^\mu)$, esto demuestra
\sref{III.6.3.6.1}.

Cuando se supone que existe $i_0$ tal que
$P_i^\lambda=Q_i^\mu=0$ para $i<i_0$, cualesquiera que sean $\lambda$ y
$\mu$, se demuestra del mismo modo que el homomorfismo canónico
\[
  \varinjlim_{\lambda,\mu}'E(P_\bullet^\lambda,Q_\bullet^\mu)
  \longrightarrow'E(P_\bullet,Q_\bullet)
  \tag{6.3.6.3}\label{III.6.3.6.3}
\]
es biyectivo.
\end{proof}

\begin{proposition}[6.3.7]
\label{III.6.3.7}
Supongamos que $P_\bullet$ y $Q_\bullet$ están acotados inferiormente. Si el
complejo $P_\bullet$ está formado por $A$-módulos planos, se tiene un
isomorfismo $A$-lineal canónico de $\partial$-funtores en $Q_\bullet$
\[
  \mathbf{Tor}_\bullet^A(P_\bullet,Q_\bullet)
  \xrightarrow{\sim}H_\bullet(P_\bullet\otimes_A Q_\bullet).
  \tag{6.3.7.1}\label{III.6.3.7.1}
\]
\end{proposition}

\begin{proof}
En efecto, la sucesión espectral \sref{III.6.3.2.1} es birregular y
degenerada, y la existencia del isomorfismo \sref{III.6.3.7.1} resulta de
\sref[0]{0.11.1.6}. Además,
calculando el hipertor a partir de una resolución proyectiva de
Cartan-Eilenberg de $P_\bullet$ como en \sref{III.6.3.4}, se ve
inmediatamente que el
isomorfismo así definido es un isomorfismo de $\partial$-funtores en
$Q_\bullet$.
\end{proof}
\begin{env}[6.3.8]
\label{III.6.3.8}
\oldpage[III]{144}
Sea $\rho:A\to A'$ un homomorfismo de anillos. Nos proponemos definir un
homomorfismo $A$-lineal funtorial de grado 0 asociado canónicamente a $\rho$:
\[
  \rho_{P_\bullet,Q_\bullet}:
  \mathbf{Tor}_\bullet^A(P_\bullet,Q_\bullet)
  \longrightarrow
  \mathbf{Tor}_\bullet^{A'}
  (P_\bullet\otimes_A A',Q_\bullet\otimes_A A').
  \tag{6.3.8.1}\label{III.6.3.8.1}
\]

Para ello, consideremos una resolución proyectiva de Cartan-Eilenberg
$L_{\bullet\bullet}$ de $P_\bullet$; consideremos, por otra parte, una
resolución proyectiva de Cartan-Eilenberg $L'_{\bullet\bullet}$ de
$P_\bullet\otimes_A A'$. Veremos que se puede definir un
homomorfismo $A'$-lineal de complejos
$L_{\bullet\bullet}\otimes_A A'\to L'_{\bullet\bullet}$, determinado salvo
homotopía. En efecto, la construcción de $L_{\bullet\bullet}$ queda
enteramente determinada cuando se elige (arbitrariamente), para cada $i$,
una resolución proyectiva $(X_{ij}^{\mathrm B})_{j\geqslant0}$ de
$\mathrm B_i(P_\bullet)$ y una resolución proyectiva
$(X_{ij}^{\mathrm H})_{j\geqslant0}$ de $\mathrm H_i(P_\bullet)$, que son
respectivamente iguales a $\mathrm B_i^{\mathrm I}(L_{\bullet\bullet})$ y
$\mathrm H_i^{\mathrm I}(L_{\bullet\bullet})$; se deduce sucesivamente
\[
  \mathrm Z_i^{\mathrm I}(L_{\bullet\bullet})
  =\mathrm H_i^{\mathrm I}(L_{\bullet\bullet})
   \oplus\mathrm B_i^{\mathrm I}(L_{\bullet\bullet}),
  \qquad
  L_{i,\bullet}=\mathrm Z_i^{\mathrm I}(L_{\bullet\bullet})
   \oplus\mathrm B_{i-1}^{\mathrm I}(L_{\bullet\bullet}).
\]
Ahora bien, $X_{i,\bullet}^{\mathrm B}\otimes_A A'$ ya no es, en general,
una resolución de $\mathrm B_i(P_\bullet)\otimes_A A'$, pero sigue siendo un
complejo formado por $A'$-módulos proyectivos, y por tanto existe un
homomorfismo $A'$-lineal
\[
  X_{i,\bullet}^{\mathrm B}\otimes_A A'
  \longrightarrow \mathrm B_i^{\mathrm I}(L'_{\bullet\bullet})
\]
compatible con las aumentaciones y determinado salvo homotopía
\emph{(M, V, 1.1)}. Se tiene asimismo un homomorfismo $A'$-lineal
\[
  X_{i,\bullet}^{\mathrm H}\otimes_A A'
  \longrightarrow \mathrm H_i^{\mathrm I}(L'_{\bullet\bullet})
\]
determinado salvo homotopía, de donde se deduce, mediante la construcción
recordada más arriba, un homomorfismo $A'$-lineal
$L_{i,\bullet}\otimes_A A'\to L'_{i,\bullet}$ para todo $i$; estos
homomorfismos (para $i\in\mathbf Z$) son compatibles con los diferenciales
$L_{i,\bullet}\to L_{i-1,\bullet}$ y con los análogos para
$L'_{\bullet\bullet}$, en virtud de la misma construcción, y constituyen por
tanto el homomorfismo $A'$-lineal
$L_{\bullet\bullet}\otimes_A A'\to L'_{\bullet\bullet}$ buscado.

Para definir \sref{III.6.3.8.1}, basta entonces considerar de igual modo una
resolución proyectiva de Cartan-Eilenberg $M_{\bullet\bullet}$ (resp.
$M'_{\bullet\bullet}$) de $Q_\bullet$ (resp.
$Q_\bullet\otimes_A A'$), y un homomorfismo $A'$-lineal
$M_{\bullet\bullet}\otimes_A A'\to M'_{\bullet\bullet}$. De estos
homomorfismos se deduce un homomorfismo $A'$-lineal
\[
  (L_{\bullet\bullet}\otimes_A A')\otimes_{A'}
  (M_{\bullet\bullet}\otimes_A A')
  \longrightarrow
  L'_{\bullet\bullet}\otimes_{A'}M'_{\bullet\bullet},
\]
y luego, por composición, un homomorfismo $A$-lineal de bicomplejos
\[
  L_{\bullet\bullet}\otimes_A M_{\bullet\bullet}
  \longrightarrow
  L'_{\bullet\bullet}\otimes_{A'}M'_{\bullet\bullet},
\]
y al pasar a la homología se obtiene \sref{III.6.3.8.1}, que está bien
definido
puesto que proviene de un morfismo de complejos determinado salvo homotopía.

Si $\rho':A'\to A''$ es un segundo homomorfismo de anillos, y
$\rho'':A\to A''$ es el homomorfismo compuesto $\rho'\circ\rho$, está claro
que
\[
  \rho''_{P_\bullet,Q_\bullet}
  =\rho'_{P'_\bullet,Q'_\bullet}\circ\rho_{P_\bullet,Q_\bullet},
  \qquad
  P'_\bullet=P_\bullet\otimes_A A',
  \quad Q'_\bullet=Q_\bullet\otimes_A A'.
\]
Observemos además que el morfismo de bicomplejos
$L_{\bullet\bullet}\otimes_A M_{\bullet\bullet}
\to L'_{\bullet\bullet}\otimes_{A'}M'_{\bullet\bullet}$ considerado
más arriba define morfismos funtoriales (en $P_\bullet$ y $Q_\bullet$) de
sucesiones espectrales
\[
  'E_{pq}^r(P_\bullet,Q_\bullet)
  \longrightarrow
  'E_{pq}^r(P_\bullet\otimes_A A',Q_\bullet\otimes_A A')
\]
y
\[
  ''E_{pq}^r(P_\bullet,Q_\bullet)
  \longrightarrow
  ''E_{pq}^r(P_\bullet\otimes_A A',Q_\bullet\otimes_A A'),
\]
independientes de las resoluciones de Cartan-Eilenberg consideradas y que
poseen también la propiedad de transitividad precedente.
\end{env}

\begin{proposition}[6.3.9]
\label{III.6.3.9}
Sea $\rho:A\to A'$ un homomorfismo de anillos tal que $A'$ sea un
$A$-módulo plano. Se tienen entonces isomorfismos canónicos funtoriales
\[
  \mathbf{Tor}_\bullet^{A'}
  (P_\bullet\otimes_A A',Q_\bullet\otimes_A A')
  \xrightarrow{\sim}
  \mathbf{Tor}_\bullet^A(P_\bullet,Q_\bullet)\otimes_A A'.
  \tag{6.3.9.1}\label{III.6.3.9.1}
\]
\[
  \left\{
  \begin{aligned}
  'E(P_\bullet\otimes_A A',Q_\bullet\otimes_A A')
    &\xrightarrow{\sim}'E(P_\bullet,Q_\bullet)\otimes_A A',\\
  ''E(P_\bullet\otimes_A A',Q_\bullet\otimes_A A')
    &\xrightarrow{\sim}''E(P_\bullet,Q_\bullet)\otimes_A A'.
  \end{aligned}
  \right.
  \tag{6.3.9.2}\label{III.6.3.9.2}
\]
\end{proposition}

\begin{proof}
\oldpage[III]{145}
En efecto, dada la exactitud del funtor $M\otimes_A A'$ en $M$,
$L_{\bullet\bullet}\otimes_A A'$ y $M_{\bullet\bullet}\otimes_A A'$ son
entonces resoluciones proyectivas de Cartan-Eilenberg de
$P_\bullet\otimes_A A'$ y $Q_\bullet\otimes_A A'$, respectivamente, de
donde resulta la conclusión.
\end{proof}

\begin{env}[6.3.10]
\label{III.6.3.10}
Sea $\rho:A\to A'$ un homomorfismo de anillos; para todo complejo
$P'_\bullet$ de $A'$-módulos, $P'_{\bullet[\rho]}$ es un complejo de
$A$-módulos; además, la aplicación identidad
$P'_{\bullet[\rho]}\to P'_\bullet$ puede considerarse como el compuesto de
las aplicaciones canónicas
\[
  P'_{\bullet[\rho]}\longrightarrow
  P'_{\bullet[\rho]}\otimes_A A'
  \xrightarrow{\mu}P'_\bullet,
\]
donde $\mu$ es el homomorfismo $A'$-lineal $\mu(x\otimes a')=a'x$. Si
$Q'_\bullet$
es un segundo complejo de $A'$-módulos, se tienen por tanto homomorfismos
canónicos funtoriales de grado 0
\[
  \mathbf{Tor}_\bullet^A(P'_{\bullet[\rho]},Q'_{\bullet[\rho]})
  \longrightarrow
  \mathbf{Tor}_\bullet^{A'}
  (P'_{\bullet[\rho]}\otimes_A A',Q'_{\bullet[\rho]}\otimes_A A')
  \longrightarrow
  \mathbf{Tor}_\bullet^{A'}(P'_\bullet,Q'_\bullet),
  \tag{6.3.10.1}\label{III.6.3.10.1}
\]
donde la primera flecha es el homomorfismo $A$-lineal definido en
\sref{III.6.3.8}, y la
segunda se deduce de los homomorfismos $A'$-lineales
$P'_{\bullet[\rho]}\otimes_A A'\to P'_\bullet$ y
$Q'_{\bullet[\rho]}\otimes_A A'\to Q'_\bullet$ por funtorialidad. Se tienen
homomorfismos análogos para las sucesiones espectrales de
\sref{III.6.3.2}, y
propiedades evidentes de transitividad, cuyo enunciado dejamos al cuidado del
lector.
\end{env}

\begin{proposition}[6.3.11]
\label{III.6.3.11}
Sea $\rho:A\to A'$ un homomorfismo de anillos que hace de $A'$ un
$A$-módulo plano. Para todo complejo $P'_\bullet$ de $A'$-módulos y todo
complejo $Q_\bullet$ de $A$-módulos acotado inferiormente, se tiene un
isomorfismo canónico funtorial
\[
  \mathbf{Tor}_\bullet^A(P'_{\bullet[\rho]},Q_\bullet)
  \xrightarrow{\sim}
  \mathbf{Tor}_\bullet^{A'}(P'_\bullet,Q_\bullet\otimes_A A').
  \tag{6.3.11.1}\label{III.6.3.11.1}
\]
\end{proposition}

\begin{proof}
En efecto, si $M_{\bullet\bullet}$ es una resolución proyectiva de
Cartan-Eilenberg de $Q_\bullet$, $M_{\bullet\bullet}\otimes_A A'$ es una
resolución proyectiva de Cartan-Eilenberg de $Q_\bullet\otimes_A A'$, y se
tiene, salvo isomorfismo canónico,
\[
  P'_{\bullet[\rho]}\otimes_A M_{\bullet\bullet}
  =P'_\bullet\otimes_{A'}(M_{\bullet\bullet}\otimes_A A');
\]
la conclusión resulta de \sref{III.6.3.4}.
\end{proof}

\begin{namedtheorem}[6.3.12]{Observación}
\label{III.6.3.12}
Sea $(A^\lambda)$ un sistema inductivo filtrante de anillos, y sean
$(P_\bullet^\lambda)$, $(Q_\bullet^\lambda)$ dos sistemas inductivos de
complejos de $(A^\lambda)$-módulos; se tiene entonces un isomorfismo canónico
que generaliza \sref{III.6.3.6.1}
\[
  \varinjlim_\lambda
  \mathbf{Tor}_\bullet^{A^\lambda}
  (P_\bullet^\lambda,Q_\bullet^\lambda)
  \xrightarrow{\sim}
  \mathbf{Tor}_\bullet^A(P_\bullet,Q_\bullet),
  \tag{6.3.12.1}\label{III.6.3.12.1}
\]
donde $A=\varinjlim A^\lambda$, $P_\bullet=\varinjlim P_\bullet^\lambda$,
$Q_\bullet=\varinjlim Q_\bullet^\lambda$. Una vez definidos, con ayuda de
\sref{III.6.3.10}, los homomorfismos
\[
  \mathbf{Tor}_\bullet^{A^\lambda}
  (P_\bullet^\lambda,Q_\bullet^\lambda)
  \longrightarrow
  \mathbf{Tor}_\bullet^{A^\mu}
  (P_\bullet^\mu,Q_\bullet^\mu)
\]
para $\lambda\leqslant\mu$, la demostración es la de \sref{III.6.3.6}.
\end{namedtheorem}

\begin{proposition}[6.3.13]
\label{III.6.3.13}
Sean $S$ un subconjunto multiplicativo de $A$, y $P_\bullet$ y $Q_\bullet$
dos complejos de $A$-módulos en los que las multiplicaciones por los
elementos de $S$ sean biyectivas, de modo que, si $A'=S^{-1}A$,
$P_\bullet$ y $Q_\bullet$ estén formados por $A'$-módulos. Entonces se tiene
un isomorfismo canónico
\[
  \mathbf{Tor}_\bullet^A(P_\bullet,Q_\bullet)
  \xrightarrow{\sim}
  \mathbf{Tor}_\bullet^{A'}(P_\bullet,Q_\bullet).
\]
\end{proposition}

\begin{proof}
En efecto, la hipótesis implica que los homomorfismos canónicos
$P_\bullet\to P_\bullet\otimes_A A'$,
$Q_\bullet\to Q_\bullet\otimes_A A'$ son biyectivos. Por otra parte, la
funtorialidad del hipertor muestra que la multiplicación por todo $s\in S$
es biyectiva en $\mathbf{Tor}_\bullet^A(P_\bullet,Q_\bullet)$ y, por
consiguiente,
\[
  \mathbf{Tor}_\bullet^A(P_\bullet,Q_\bullet)
  \longrightarrow
  \mathbf{Tor}_\bullet^A(P_\bullet,Q_\bullet)\otimes_A A'
\]

\oldpage[III]{146}

es también un homomorfismo biyectivo. Como $A'$ es un $A$-módulo plano, la
conclusión resulta de \sref{III.6.3.9}, y se tienen asimismo isomorfismos
canónicos
para las sucesiones espectrales.
\end{proof}
